\section{Related Work}
\label{sec:lit}

\subsection{Streaming and Treatment-Area Design}

Streaming separates patients into flows that are treated in different areas or by different teams. The most common form is a fast track for patients with less severe conditions. A systematic review of triage-related interventions found moderately strong evidence that fast tracks reduce waiting time, LOS and the share of patients who leave without being seen, but limited evidence for other forms of streaming \cite{Oredsson2011_TriageFlowReview}. Before--after studies of streaming in individual EDs report improvements for low-acuity patients \cite{Kinsman2008,King2006} and mixed effects on other outcomes \cite{StreamingINED_2015_Kim_etal}; a realist analysis argues that the effect depends on local mechanisms and context \cite{StreaminINED}. Operations research has examined when streaming helps. Saghafian et al.\ \cite{Saghafian2012_Streaming_as_a_Mechanism} studied streaming by predicted disposition and showed that its benefit depends on the operating conditions. Queueing results on partitioning servers and customers show that separating patients with different service needs can reduce waiting, but at the cost of lost pooling \cite{Whitt1999,Hu2009}. Flexible partitioning between low- and high-acuity areas has been proposed to recover part of that loss \cite{laker2014flex_track}. Split or vertical flow, in which ambulatory patients are assessed in chairs and do not hold a bed while waiting for results, has been evaluated by simulation in single EDs \cite{Konrad2013_SplitFlowSim} and analysed as a streaming decision \cite{feizi2023vertical}.

\subsection{Prioritization Rules}

Triage scales such as the CTAS assign each patient an acuity level and a target time to physician assessment \cite{Bullard2017_CTAS}. Strict priority by triage level protects the most urgent patients but can leave low-acuity patients waiting for very long times. Empirical work shows that actual prioritization in Canadian EDs is more nuanced than strict CTAS order \cite{Prioritization_Using_CTAS}. The accumulating priority queue (APQ) lets priority grow with waiting time at a class-specific rate, so that long waits are bounded more evenly across classes \cite{StanfordTaylorZiedins2014_APQ}; ED simulation studies report that APQ-type assignment can reduce LOS \cite{Vanbrabant2021}. Other work considers prioritization under ED blocking \cite{Li_etal_2023-Who-Is-Next}, dynamic patient-to-physician assignment \cite{Dynamic_scheduling_ED2023EJOR}, the use of complexity information at triage \cite{Saghafian2014_Complexity-Augmented_Triage}, provider triage \cite{Provider_Triage_Management_Science_2019} and fair patient placement \cite{POM_Patient_PlacementInED_2024}.

\subsection{Prediction-Based Dispatching}

Machine learning models can predict admission, resource use and LOS at triage \cite{Hong2018_disposition,SUN2011_disposition,Azari2015_LongLOS}. Whether such predictions improve operations depends on how they are used and on their accuracy. Studies have characterized the value of predictive analytics for patient flow \cite{Peck_2014_value_of_predictive_analytics}, examined when ED decisions should use predictions at all \cite{POM_To_Predict_or_Not_2022}, and compared queueing and machine learning approaches to delay prediction \cite{Chocron2024DelayPrediction}. A broader literature studies how predictions should enter prescriptive decisions \cite{BertsimasKallus2020_Prescriptive,ElmachtoubGrigas2022_SPO}. For ED streaming, a practical question is the accuracy threshold below which a prediction-based policy no longer beats a policy that uses triage information only.

\subsection{Similarity-Based Rules}

Similarity coefficients were developed in group technology to compare coded parts and form part families \cite{OFFODILE1990,Seifoddini01091986,YIN2005471}. They have been carried over to EDs to group and prioritize patients with similar clinical and operational profiles \cite{GT-in-ED_HCMS2016}. A similarity rule ranks patients by the weighted closeness of their attribute code to a target code. The appeal for managers is that one representation can express many policies by changing weights and targets. The same flexibility raises a question that, to our knowledge, has not been examined: for which weights and targets does a similarity rule implement a given classical rule, and how far is it from that rule otherwise?

\subsection{ED Simulation Across Settings and Performance Measures}

Discrete-event simulation is the standard tool for evaluating ED operations \cite{Gul2015_SimInED,Saghafian2015}. Reviews of ED simulation list a wide range of performance indicators and interventions, mostly evaluated in one ED at a time \cite{VANBRABANT2019356_Simulation_of_ED}. Performance measurement in EDs relies on time intervals such as door-to-physician time and LOS \cite{Sorup2013_EDPerformance}, with standard time-stamp definitions \cite{ED-Metrics_WELCH2013}; crowding arises from input, throughput and output factors, including boarding of admitted patients \cite{Asplin2003EDCrowding,hoot2008systematic}. Simulation-based metamodels have been used to support ED resource planning \cite{ZEINALI2015123}. Reporting guidelines for simulation studies ask for transparent description of logic, data, experimentation and implementation \cite{Monks2019_STRESS}.

\subsection{Gap}

The literature offers many streaming and prioritization policies, each evaluated in one or a few settings, and the conclusions do not always agree. We did not find a study that compares the same set of streaming policies across a designed range of ED types, measures the loss from applying one policy to all of them, and checks whether scoring-rule implementations apply the policies they are meant to represent. This paper addresses that gap.
